\scanpage{164}
\deleted{Lorsque les} Bien sur, les différentielles \inserted{(6)} définissent
\[
 d_1^{p,q}[p]:\mathscr H^q\bigl(\Omega^p_{S/T}\otimes\mathbb R f_*(\Omega^\bullet_{X/S})\bigr)\longrightarrow
 \mathscr H^q\bigl(\Omega^{p+1}_{S/T}\otimes\mathbb R f_*(\Omega^\bullet_{X/S})\bigr) \,, \tag*{(\deleted{\unsure{6}}8)}
\]
qui, pour $\Omega^p_{S/T}$ plat sur $S$, p. ex $S$ formellement lisse sur $T$, \deleted{\unsure{def}} s'écrit aussi
\[
 d_1^{p,q}[p]:\Omega^p_{S/T}\otimes\mathbb R^q f_*(\Omega^\bullet_{X/S})\longrightarrow
 \Omega^{p+1}_{S/T}\otimes\mathbb R^q f_*(\Omega^\bullet_{X/S}) \,, \tag*{(8$'$)}
\]
qui pour $p=0$ définit une connexion intégrable sur les $\mathbb R^q f_*(\Omega^\bullet_{X/S})$, dite conn. de \uline{Gauss-Manin}.

Posons pour un moment, pour simplifier,
\[
 K^\bullet=\mathbb R f_*(\Omega^\bullet_{X/T}) \,, \tag*{(9)}
\]
qu'on considère comme complexe filtré comme on l'a dit. On aura donc
\[
 \mathbb R\Gamma_S(K^{\unsure{\bullet|'}})\simeq\mathbb R\Gamma_X(\Omega^\bullet_{X/T})\simeq\mathbb RH_{\mathrm{DR}}(X/T). \tag*{(10)}
\]
La filtration de $K^\bullet$ permet, d'après un \unsure{sorte|sorite|sortie} général, d'écrire une suite spectrale
\[
 \mathbb R\Gamma_S(K^\bullet)\Longleftarrow E^{\unsure{p,q|pq}}_2=\mathbb H^p(S,\alpha\longmapsto\mathscr H^{\alpha+q}(\mathrm{Gr}^\alpha(K^\bullet))) \,; \tag*{(11)}
\]
les différentielles définissant les complexes $\alpha\mathrel{\unsure{\longmapsto|\mapstochar\relbar\mapstochar}}\ldots$ étant la différentielle évidente $d_1$ provenant de la filtration. Tenant compte de (3) \inserted{et (10)}, on trouve donc la suite spectrale.
% archive sheet number: 163
